\section{Introduction}

Anticipating hand motion is a useful capability in egocentric and interactive settings, where the hand is both the primary actuator \cite{GarciaHernando2018FirstPerson} and a strong cue for upcoming action \cite{Liu2022Hotspots,Hatano2025EgoH4}. In embodied assistance, teleoperation, and human--robot interaction, forecasting how the hand will move in the near future can inform intended manipulation and support anticipatory interaction \cite{Kothari2023ConstrainedHRCMotionPrediction,Heman2026TeleoperationIntent}. Such forecasts can also provide motion cues that may be useful in assistive monitoring settings. More broadly, the goal is not only to recognize what a user is doing now, but also to model how the motion is likely to unfold over the next few steps.

This forecasting problem is especially demanding for 3D hand motion. Compared with broader articulated motion, hand trajectories are more densely structured, more locally coupled, and more sensitive to fine-grained deformation associated with grasping or interaction. The future therefore contains both stable trend-like components and ambiguous residual variations: some aspects of the motion are predictable from context and recent dynamics, while others remain multi-modal at the level of individual articulators. A practical model must capture both aspects without sacrificing articulated consistency.

Existing approaches only partially address this combination of structure and uncertainty. Deterministic sequence predictors can capture dominant motion trends effectively \cite{Martinez2017RNN,Mao2020HistoryRepeats,Dang2021MSRGCN}, and strong feed-forward baselines further reinforce this point \cite{Guo2023SiMLPe}, but such models tend to collapse multiple plausible futures into a single average prediction. Stochastic formulations can represent variability more explicitly \cite{Yuan2020DLow,aliakbarian2020stochastic,Chen2023HumanMAC}. More recent diffusion-based predictors pursue the same goal from a generative perspective \cite{Sun2024CoMusion,Tang2024PromptFDDM}, yet generic generative formulations do not necessarily encode articulated hand structure or the fact that future uncertainty may be uneven across motion frequencies. This motivates a forecasting formulation that separates predictable trajectory trend from residual uncertainty while structuring the latter around hand-specific motion geometry.

We address this problem with \modelname{}, a coarse-to-fine model for 3D hand-motion forecasting. The method first predicts a deterministic coarse future in the frequency domain, using a compact DCT representation to capture the dominant spectral structure of the trajectory. It then refines this prediction with a conditional diffusion model defined over future-frequency residuals. The residual process is organized by a spatio-spectral prior that combines articulated hand topology with frequency-dependent future modes, so that uncertainty is modeled along directions induced by articulated hand structure and future-frequency variation. In this way, \modelname{} treats forecasting as the combination of coarse deterministic prediction and structured stochastic residual refinement. The full formulation is detailed in Sec.~\ref{sec:method}.

Our contributions are as follows:
\begin{itemize}
    \item We introduce a frequency-domain coarse-to-residual formulation for 3D hand-motion forecasting, in which a deterministic predictor captures the dominant future trajectory and a stochastic model focuses on the remaining residual uncertainty.
    \item We develop a conditional diffusion model over future-frequency residuals, where sampling refines a coarse forecast rather than generating the entire trajectory from scratch.
    \item We define a spatio-spectral anisotropic residual prior that combines articulated hand topology and future-frequency modes to structure the residual uncertainty space.
    \item We evaluate the resulting model against stochastic and deterministic forecasting baselines on standard 3D hand-motion benchmarks, with quantitative results and ablations reported in Sec.~\ref{sec:experiments}.
\end{itemize}
